% !TEX root = ../CINTA-cn.tex
%%%%%%%%%%%%%%%%%%%%%%%%%%%%%%%%%%%%%%%%%%%%%%%%%%%%%%%%%%%%%%%%%%%%
% Author: Libin Wang
% Chapter 12. Ring and Field
% This document was released as open source on 2026-06-18.
% Copyright 2023, Libin Wang. Creative Commons License
% This work is licensed under a Creative Commons Attribution-NonCommercial-NoDerivatives 4.0 International License.
%%%%%%%%%%%%%%%%%%%%%%%%%%%%%%%%%%%%%%%%%%%%%%%%%%%%%%%%%%%%%%%%%%%%

\chapter{环与域}
\label{chap:ring_field} 
\begin{introduction}
  \item 环
  \item 整环
  \item 理想
  \item 多项式环
  \item 域
  \item 有限域
\end{introduction}
代数结构\emph{群}针对一个特定的集合，定义该集合中元素的一种操作，不是“加法”就是“乘法”，
比如整数的加法群$\Z_n$和整数的乘法群$\Z_n^*$。
然而常识告诉我们，对于整数应该是同时具有加法与乘法，为什么不定义一种同时具有两种元素操作的代数结构呢？
这一章讨论两种代数结构：\emph{环}（\emph{ring}）\index{环}\index{ring}和
\emph{域}（\emph{field}）\index{域}\index{field}，它们都定义两种元素操作。

\section{环}
\label{sec:ring}
在计算机系统中，最自然的代数结构就是环。
环是一种包括两种二元操作的代数结构，具体定义如下：

\begin{definition}{环}{ring}
环$R$是一个非空集合，在$R$上有两种封闭的二元操作：加法（记为$+\colon R \times R \to R$）
和乘法（记为$\cdot\colon R \times R \to R$），并且满足以下条件：
\begin{enumerate}
\item 在加法（$+$）上$R$是一个阿贝尔群，加法的单位元记为$0$，加法上的逆元记为$-a$；
\item $R$在乘法（$\cdot$）上满足结合律；
\item 乘法在加法上满足分配律。
\end{enumerate}
\end{definition}
具体地表示为以下公式，对任意$a,b, c\in R$，环$R$满足以下公理：
\begin{align}
a + (b + c ) &= (a + b) + c   \quad \quad \text{($+$的结合律)}\\
a + b  &=  b + a   \quad \quad \text{($+$的交换律)}\\
a + 0 &= 0 + a \quad \text{($+$的单位元)}\\
a + (- a) &= (- a) + a = 0 \quad \quad \text{($+$的逆元)}\\
a \cdot (b \cdot c) &= (a \cdot b) \cdot c \quad  \quad\text{($\cdot$的结合律)}\\
% a \cdot 1= 1\cdot a = a \quad \quad \text{($\cdot$的单位元)}\\
(a + b) \cdot c &= (a \cdot c) + (b \cdot c) \quad \quad \text{(右分配律)}\\
a \cdot (b + c) &= (a \cdot b) + (a \cdot c) \quad \quad \text{(左分配律)}
\end{align}

如果$R$在乘法上也满足交换律，即对任意的$a, b \in R$，有
\[ a \cdot b = b \cdot a \]
则称$R$为\emph{交换环}，否则称为非交换环。
以后，在不发生歧义时，表述环的乘法往往省略乘法符号，
比如把$a \cdot b$写成$ab$。

如果$R$在乘法上具有单位元，即存在$1 \in R$，$1 \neq 0$，且对任意的$a \in R$，
有$a 1 = 1 a = a $，则称环$R$为带单位元的环。以下我们重点讨论带单位元的环，不特别声明
都假定环带单位元。请看以下若干环的实例，建立对环最直观的认识。
\begin{example}
\begin{enumerate}
\item $R = \{0\}$是只有一个元素的最小环，称为平凡环或零环。如果一个环中，$1 \neq 0$，则这个环是非平凡环。
一个最小的非平凡环就是$R = \{0, 1\}$，或者记为$\Z_2$。

\item 在普通意义的加法和乘法上，容易验证$\Z, \mathbb{Q}, \mathbb{R}, \mathbb{C}$都是交换环。比如，考虑$\Z$，它
在加法上是群，在乘法上满足结合律，有单位元$1$，交换律和分配律则是大家熟知的事实。

\item 整数中所有的偶数在一般的加法与乘法上形成环，同样记为$2\Z$，这是一个不带单位元的环。
           更一般地，对任意整数$n\in \Z$，$n\Z$在一般的加法与乘法上形成环。

\item 对任意正整数$n\in \N$，$\Z_n$是模$n$的加法群。如果给加法群$\Z_n$配置上乘法$\cdot$，
乘法$\cdot$被定义为整数上的模$n$乘法，即任取$a, b \in \Z_n$，
$a \cdot b \triangleq ab\bmod n$。容易验证$\Z_n$是一个交换环。
注意，此时$\Z_n$的元素在乘法上不一定存在逆元。$\Z_n$是计算机系统中最基本的代数结构，
因为计算机系统中的算术都是模$n$的加法和模$n$的乘法，并且要求$n$是$2$的某个指数。

\item 不难想到，如果取$p$为任意素数，$\Z_p$在模$p$的加法与模$p$的乘法下成环，
而且$\Z_p$的所有非零元素在乘法上都有逆元。以后我们将重点讨论这种特殊的环。

\item 对任意环$R$，$R$的直积$R \times R$（参考群论的定义~\ref{def:dir_prod}）形成环，环的加法与乘法定义
为：对任意的序对$(a,b), (c,d) \in R \times R$，
$(a, b) + (c, d) \triangleq (a+c, b+d)$和$(a, b)  (c, d) \triangleq (ac, bd)$。

\item 实数上的$n \times n$矩阵在普通的矩阵加法和矩阵乘法上形成环，也称为矩阵环，记为$M_n(\mathbb{R})$。这是一种非交换环。
\end{enumerate}
\end{example}

%%%20201229
环定义中要求加法群满足交换律，初看上去会感觉这是一种人为的强要求，其实不然。因为，如果环中包含乘法单位元，则加法交换律必然成立。
设$R$是环，任取$a, b \in R$，考虑$(a + b)(1 + 1) $，分别应用左分配律和右分配律，有：
\[(a + b)(1 + 1) = (a + b) + (a + b) = (a + a) + (b + b)  \]
所以有$a + b = b + a$，即加法交换律成立。

\begin{proposition}{}{ring0}
设$R$是一个环，且$a, b \in R$，则有：
\begin{enumerate}
\item a0 = 0a = 0
\item a(-b) =  (-a)b = -ab
\item(-a)(-b) = ab
\end{enumerate}
\end{proposition}

\begin{proof}
根据分配律，
\[ a 0 = a  (0 + 0) = a0 + a0\]
则$a0 = 0$，同理$0a = 0$。同样根据分配律
\[ab + a(-b) = a(b + (-b)) = a 0 = 0\]
所以，$-(ab) = a(-b)$，同理$-(ab) = (-a)b$。最后，
根据以上结论，
$(-a)(-b) = -(a(-b)) = -(-(ab)) = ab$。
\end{proof}

\begin{remark}
命题~\ref{pro:ring0}中的性质似乎非常平凡，考虑环$\Z$或$\Z_n$，这些都是我们熟知的性质。
然而，考虑到抽象代数的“抽象性”，则不可想当然地认为这些属性必然成立。
\end{remark}

环并不要求环元素在乘法上存在逆元，一个直接的结论就是：在乘法上，消去律不必然存在！
如何修补定义而获得乘法的消去律呢？最自然的想法就是让环中所有非零元素在乘法上成群。进而再问，
这是否必须的呢？以上问题引出了\emph{整环}（\emph{integral domain}）\index{整环}\index{integral domain}的概念。
\begin{definition}{整环}{int_domain}
给定一个环$R$，对非零元素$a \in R$，如果存在非零元素$b\in R$，且$a \cdot b = 0$，
则称$a$为一个\emph{零因子}（\emph{zero divisor}）。如果交换环$R$中没有零因子，
即任取$a, b \in R$，如果$ab = 0$则有$a = 0$或$b = 0$， 
则称$R$为整环。
\end{definition}

\begin{example}
\begin{enumerate}
\item 在普通意义的加法和乘法上，$\Z, \mathbb{Q}, \mathbb{R}, \mathbb{C}$都是整环。比如，考虑$a, b \in \Z$，
$ab = 0$ 当且仅当$a=0$或者$b = 0$。其他实例，验证类似。

\item 可验证偶数环$2\Z$是整环。更一般地，对任意的$n\in \Z$，$n\Z$都是整环。

\item 已知$\Z_n$是一个交换环。但是，$\Z_n$并不是整环。比如，$\Z_{15}$中，$(3 \cdot 5) \bmod 15 = 0$，
$3$和$5$都是$\Z_{15}$中的零因子。

\item 矩阵环$M_n(\mathbb{R})$不是整环，因为存在$A, B \in M_n(\mathbb{R})$，使得$AB = 0$但是$A$
和$B$都不为$0$。
\end{enumerate}
\end{example}

\begin{proposition}{整环的消去律}{ring_cancel}
设$D$是一个交换环，$D$是整环当且仅当对任意元素$a,b,c\in D$，且$a \neq 0$，若$ab = ac$，则$b = c$。
\end{proposition}

\begin{proof}
\begin{enumerate}
\item $\Rightarrow$. $D$是整环，则$D$中无零因子。若$a \neq 0$且$ab = ac$，则$a(b - c) = 0$，则$b - c = 0$，即$b = c$。

\item $\Leftarrow$. 假设$D$中消去律成立。任取$a, b \in D$且$a \neq 0$。
设$ab = 0$，则有$ab = a0$。根据消去律，
得到$b = 0$。因此，$a$不可能是零因子。
\end{enumerate}
\end{proof}
命题~\ref{pro:ring_cancel}解答了之前提出的问题，整环可确保在较弱的条件下环的乘法消去律存在。
有一点值得强调，根据定义要求，整环必然是交换环，通过下一节的某个证明我们可以更深刻地体会这个要求的必要性。%20230619增加

\begin{definition}{子环}{subring} \index{子环}\index{subring}
给定环$R$，$R'$是$R$的子集，如果$R'$在环$R$的加法和乘法上也形成环，则称$R'$是$R$的子环，记为$R' \subset R$。
\end{definition}

\begin{example}
\begin{enumerate}
\item 容易验证以下子环序列：$\Z\subset \mathbb{Q}\subset \mathbb{R}\subset \mathbb{C}$。

\item 偶数环是整数环的子环，即$2\Z \subset \Z$。由此可知，子环并不自然继承母环的单位元。
            当然，对任意的$n\in \Z$，有$n \Z \subset \Z$。
\end{enumerate}
\end{example}

\begin{proposition}{}{subring}
给定环$R$，$R'$是$R$的子集。$R'$是$R$的子环，当且仅当以下条件满足：
\begin{enumerate}
\item $R' \neq \emptyset$；
\item $\forall a, b \in R'$，有$ab \in R'$；
\item $\forall a, b \in R'$，有$a - b \in R'$。
\end{enumerate}
\end{proposition}

\begin{proof}
根据命题~\ref{pro:subgr}易得，留作课后练习。
\end{proof}

第~\ref{chap:iso_hom}章已经建立了一条紧密结合的知识链：同构、同态、kernel、正规子群、商群和同构定理。
同样的，在环论的知识体系中也存在类似的对应物，也包括：同构、同态、kernel、商环、同构定理等等。
以下简要给出相关定义与定理，并重点讨论在环的语境下与群论中正规子群对应的一个概念：\emph{理想}（\emph{ideal}）\index{理想}\index{ideal}。
\begin{definition}{环同态与环同构}{ring_iso_homo}
给定两个环$R$和$R'$，若映射$\phi\colon R \to R'$满足：$\forall a, b \in R$，
\begin{align*}
\phi(a + b) &= \phi(a) + \phi(b)\\
\phi(ab) &= \phi(a)\phi(b)
\end{align*}
则称$\phi$为环同态。如果$\phi$是一种双射，则$\phi$是环同构。
环同态$\phi$的 kernel 定义为以下集合：
\[\Ker\; \phi = \{r \in R: \phi(r) = 0\} \]
上式中的$0$指的是环$R'$的加法单位元。
\end{definition}

\begin{proposition}{环同态的性质}{ring_homo}
设映射$\phi\colon R \to R'$是环同态，则：
\begin{enumerate}
\item 如果$R$是交换环，则$\phi(R)$也是交换环。

\item 分别记$0$和$0'$是$R$和$R'$的加法单位元，$\phi(0) = 0'$。

\item 分别记$1$和$1'$是$R$和$R'$的乘法单位元，如果$\phi$是满射，则$\phi(1) = 1'$。
\end{enumerate}
\end{proposition}

\begin{proof}
留作课后练习。
\end{proof}

%examples of ring homomorphisms
%https://math.stackexchange.com/questions/1241280/if-r-is-a-ring-with-unit-element-1-and-varphi-is-a-homomorphism-of-r-on
% 20200711
\begin{example}
\label{exm:homo_onto}
命题~\ref{pro:ring_homo}的证明与群论中对应的证明类似，依葫芦画瓢很容易得到证明。必须提醒大家要思考的是，
命题为什么要求$\phi$是满射才会有$\phi(1) = 1'$？请验证以下同态实例，体会环同态与群同态的异同点。

\begin{enumerate}
\item 已知$\Z$是环，定义映射$\phi\colon \Z \to \Z$为$\phi(k) = 2k$，$\forall k \in \Z$，即把所有整数映射到偶数$2\Z$。
已知$\phi$是$\Z$到$\Z$的群同态，但是可以验证$\phi$\emph{不是}一种环同态。

\item 对任意整数$n\in \Z$，已知$\Z_n$为环，定义映射$\phi\colon \Z_n \to \Z_n$为$\phi(a) = a^2 \bmod n$，$\forall a \in \Z_n$，
即把所有$\Z_n$的元素映射到平方数。可以验证$\phi$\emph{不是}一种环同态，尽管这种映射在$\Z_n^*$中是一种群同态。

\item 考虑一种特殊的环同态，定义$\phi\colon 2\Z \to \Z_2$为$\forall k \in 2\Z$，$\phi(k) = k \bmod 2$。明显，$\phi$
是同态，但不是满同态，它把所有偶数都映射到了$0$。将这种把任意$R$映射到零环$\{0\}$的环同态称为\emph{零同态}。

\item 以下是一种更特殊的环同态，对任意的环$R$，定义映射$\phi\colon R \to \Z_2$为$\forall a \in R$，$\phi(a) = 0$。可验证，
这是一种零同态，显然不是满同态。特别提醒注意，如果环$R$有乘法单位元$1$，$\phi(1) = 0$。并不是我们期望的映射到$\Z_2$的$1$。

\item 对任意环$R$，已知$R\times R$是环。定义映射$\phi\colon R\times R \to R\times R$为
$\forall (a, b)\in R\times R$，$\phi(a,b) = (a, 0)$。可验证，这是一个环同态，但并非满同态。注意，
在$\phi(R\times R)$中的单位元是$(1,0)$，但是$R\times R$中的单位元是$(1,1)$。

\item 设$p$是任意一个奇素数，考虑$2\Z$与$\Z_p$之间的映射$\phi\colon 2\Z \to \Z_p$，
定义为$\forall k \in 2\Z$，$\phi(k) = k \bmod p$。
直观上看，该映射把所有的偶数做模$p$操作满射到环$\Z_p$。容易验证$\phi$是满同态。值得注意的是，
偶数环$2\Z$没有乘法单位元，环$\Z_p$的乘法单位元是$1$，$2\Z$中有无穷多的元素映射到$\Z_p$的乘法单位元$1$上。
也就是说，即使乘法单位元必然映射为乘法单位元，也并非只有乘法单位元才映射为乘法单位元。

\end{enumerate}
\end{example}

环同态的定义显得很自然直接，与群同态的定义、属性也非常类似，但是，只要揣摩一下以上实例还是能明显感觉到
环同态与群同态的重大区别。接下来介绍的这个定义则与群论的对应物有明显区别。
回顾群论中的正规子群，它的性质特殊，与 kernel 有密切关系，并用于定义商群，有着非常重要的作用。
在环论中与之对应的概念是\emph{理想}。

回顾群论中的定义，子群$\mathbb{H}$是群$\mathbb{G}$的正规子群，
则满足对任意的$g \in \mathbb{G}$，
$g^{-1} \mathbb{H} g = \mathbb{H}$；
或者等价地，要求
$g \mathbb{H}  = \mathbb{H} g$。与之不同，
由于环的定义并不要求元素在乘法上有逆元，所以，这就决定了环中的理想有着与正规子群不同的形式。

\begin{definition}{理想}{ideal}\index{理想}\index{ideal}
给定环$R$，$I$是$R$在加法上的子群。如果对任意的$r \in R$有$rI \subset I$和$Ir \subset I$，则称$I$是$R$的理想。如果$I$是$R$的理想，又是$R$的真子集，
则称$I$是$R$的真理想（proper ideal）\index{真理想}\index{proper ideal}。
\end{definition}
从表面上看，所谓环$R$的理想$I$，首先它是$R$的加法群的子群，其次它具有“\emph{吸收性}”，即对任意的环元素$r\in R$，无论它是否落在$I$中，
用$r$左乘或者右乘$I$，所得到的元素都会落回到$I$中。形式化的表达就是，对任意的$r\in R$和任意的$i \in I$，
必然有$ri\in I$和$ir \in I$。在此，需要读者去体会理想的“吸收性”与正规子群的“左陪集与右陪集相等”的异同之处。

\begin{remark}
\label{rem:ideal}
很多教科书是把理想定义为环的子环，这是一种方便的做法。
其实，因为理想具有“吸收性”，可容易验证，理想在环的乘法上满足封闭性，加上理想是环加法上的子群，所以，理想也是环的子环。
\end{remark}

\begin{example}
\label{exm:ideal1}
\begin{enumerate}
\item 所有的环$R$都有两个平凡理想：$\{0\}$和$R$。

\item 如果$R$的理想$I$中包括$1$，则$R=I$。请读者自己证明。

\item 对任意整数$n\in \Z$，集合$n\Z$是环$\Z$的理想。直观上看，集合$n\Z$包含了所有$n$的倍数，$n\Z$在加法上成群，
而用任意整数乘$n$的倍数还是得到一个$n$的倍数，虽然此时$n\Z$中并不必然有单位元$1$，也不必然有乘法逆元。

\item 如果$I$和$J$是环的理想，那么$I + J = \{i + j : i \in I, j \in J\}$也是理想。留作课后练习。
\end{enumerate}
\end{example}

%%%
%20201209
%AATA
%https://math.stackexchange.com/questions/1102037/the-chinese-remainder-theorem-for-rings
%%%
利用理想，可以定义环中元素\emph{模}理想的\emph{同余关系}。设$R$是环，$I$是$R$中的理想，那么对任意的$a, b \in R$，
如果$a \in b + I$，则称$a$与$b$模$I$同余，并记为：
\[ a \equiv b \pmod{I}\]
或者说，$a$与$b$模$I$同余意味着，存在$i \in I$使得$a = b + i$，等价地即$a - b \in I$。%%%利用群论的语言，也可以说a和b落在相同的陪集。

\begin{example}
已知，$\Z$是环，$2\Z \subset \Z$是$\Z$的理想。那么
\[ 7 \equiv 5 \pmod{2\Z}\]
因为，$7 = 5 + 2$，且$2 \in 2\Z$。但是，
\[ 7 \not\equiv 6 \pmod{2\Z}\]
因为，不存在偶数加$6$会等于$7$。
\end{example}

\begin{thinking}
\label{rem:mod_ideal}
请思考，模理想的同余关系与模整数的同余之间是什么关系呢？提示：给定整数$n$，$n\Z$是$\Z$理想，$a$与$b$模$n$同余意味着存在某个整数$k$使得$a - b = kn$。
用理想来表示，不正是$a - b \in n\Z$吗？
\end{thinking}

利用模理想的同余关系，可以非常容易地得到中国剩余定理在环中的版本。
\begin{theorem}{中国剩余定理--环版本}{crt_ring}
设$R$是环，$I$和$J$是$R$中的理想，并且$I + J = R$。 对任意$r_1, r_2 \in R$，以下方程组有解，并且，该方程组的任意解都模$I \cap J$同余。
\begin{align*}
x &\equiv r_1 \pmod{I}\\
x &\equiv r_2 \pmod{J}
\end{align*}

\end{theorem}

\begin{proof}
因为$I + J = R$，所以存在$i \in I$和$j \in J$使得$i + j = r_2 - r_1$。注意，此时只需要把$r_2 - r_1$理解为$R$中的某个元素即可。
%%取r_2 - r_1是为了方便凑出以下形式！
令$x' = r_1 + i = r_2 - j$，可知$x' \in r_1 + I$且$x' \in r_2 + J$，所以$x'$是方程组的解。

假设方程有两个解$x_1$和$x_2$，那么必然有：
\begin{align*}
x_1 &\equiv x_2 \pmod{I}\\
x_1 &\equiv x_2 \pmod{J}
\end{align*}
则有$x_1 - x_2 \in I$和$x_1 - x_2 \in J$，因此$x_1 - x_2 \in I \cap J$，即
\[ x_1 \equiv x_2 \pmod{I \cap J}\]
\qed
\end{proof}

\begin{example}
已知，$\Z$是环，$2\Z, 3\Z \subset \Z$是$\Z$的理想，且$2\Z + 3\Z = \Z$。求解：
\begin{align*}
x &\equiv 5 \pmod{2\Z}\\
x &\equiv 4 \pmod{3\Z}
\end{align*}
易知：$ 7 \equiv 5 \pmod{2\Z}$，且$7 \equiv 4 \pmod{3\Z}$。
$7$是唯一解吗？答案是否定的，因为$2\Z \cap 3\Z = 6\Z$，所以满足以下同余式的$x$都是解：
\[ x \equiv 7 \pmod{6 \Z}\]
\end{example}
%20200706

%20241204下午，根据课后学生的提问，产生了灵感。
\begin{thinking}
\label{rem:ideal_crt}
请思考，环版本的 CRT 要求作为“模数”的两个理想$I$和$J$满足$I + J = R$，数论版本的 CRT 要求各同余式的模数之间两两互素，请问这两者之间存在何种关系？
提示：给定任意两个整数$n_1$和$n_2$，$n_1\Z+ n_2\Z$什么条件下会等于$\Z$？进一步，环版本的 CRT 是在模$I\cap J$下有唯一解，
数论版本的 CRT 在是模各个同余式的模数的乘积之下有唯一解，请问这两者之间又存在何种关系？
%%gcd(n_1, n_2) = 1 iff $n_1\Z+ n_2\Z =\Z$
\end{thinking}

\begin{proposition}{主理想}{prin_ideal}
设$R$是一个交换环且有单位元，任取$a \in R$，定义集合
\[ \langle a \rangle \triangleq \{ar: r \in R\}\]
可验证$\langle a \rangle$是环$R$的一个理想，并称之为\emph{主理想}（\emph{principal ideal}）\index{主理想}\index{principal ideal}。
\end{proposition}

\begin{proof}
首先，验证$\langle a \rangle $是非空集合，至少包括$0$和$a$两个元素。
其次，验证$\langle a \rangle$在加法上成群。
最后，验证$\langle a \rangle$ 具有吸收性，即任取$s \in R$乘上$\langle a \rangle $中任意元素$ar$，
必然有
\[s (ar) =  a(sr) \in \langle a \rangle \]
注意，上式成立需要依赖交换律。所以，$\langle a \rangle$是$R$的理想。
\end{proof}

\begin{example}
\label{exm:ideal2}
\begin{enumerate}
\item 对任意环$R$，只包含一个元素$0$的主理想$\langle 0 \rangle$称为\emph{零理想}。

\item 对任意带单位元的环$R$，称$\langle 1 \rangle$为\emph{单位理想}（\emph{unit ideal}），显然$R = \langle 1 \rangle$。
\index{单位理想}\index{unit ideal}

\item 对任意整数$n$，集合$n\Z$是整数环$\Z$的理想，也是主理想，$n\Z = \langle n \rangle$。
\end{enumerate}
 
\end{example}

\begin{proposition}{}{int_prin_ideal}
整数环$\Z$的所有理想都是主理想。
\end{proposition}

\begin{proof}
首先，零理想也是主理想，因为$\langle 0 \rangle = \{0\}$。设$I$是整数环$\Z$的一个非零理想，
则$I$中必然包括某些正整数，根据良序原则（公理~\ref{axi:POWO}），
则$I$中必然存在一个最小正整数$n$。对任意的元素$a\in I$，根据除法算法，
存在整数$q$和$r$，$0 \leq r < n$，使得：
\[ a = qn + r\]
也就是，$r = a - qn$，利用理想的属性，可知$r \in I$。又因为$n$是$I$中最小正整数，
所以，$r=0$。因此，$a = qn$，即$I = \langle n \rangle $。
\end{proof}

环论中的主理想对应于群论中的循环群，循环群$\langle g \rangle $是包含$g$的最小群，
而理想$\langle a \rangle$则是包含$a$的最小理想。命题~\ref{pro:sub_cyc}告诉我们，
循环群$\mathbb{G}=\langle g \rangle$的每一个子群都是循环群，但是，针对主理想，则没有类似的结论。
即如果$I$是$R$的主理想，$I'$是$I$的子集，且是$R$的理想，则$I'$不一定是主理想。

\begin{proposition}{}{ring_kernel}
环同态$\phi\colon R \to R'$的 kernel 是$R$的理想。
\end{proposition}

\begin{proof}
根据群论的结论，$K = \Ker\;\phi$是$R$的加法子群（并且是正规子群）。
那么，只需要证明$K$具有理想的“吸收性”，
即对任意的$r \in R$和$a \in K$有$ar \in K$和$ra \in K$。
显然如此，因为：
\[\phi(ar) =  \phi(a) \phi(r) = 0 \phi(r) = 0 \]
且
\[\phi(ra) = \phi(r )\phi(a) = \phi(r) 0 = 0\]
\end{proof}

\begin{example}
对任意整数$n\in \Z$，定义映射$\phi\colon \Z \to \Z_n$为：对任意$a \in \Z$，$\phi(a) = a \bmod n$。容易验证
这是一个环同态，而$\Ker\;\phi$就是$n\Z$。
\end{example}
%20200710
讲完理想，是时候开始讨论商环了。在证明商环定理之前，让我们先回顾、理解、熟悉相关的思路。首先，理解商环必须
从商群出发。因为一个环$R$，本身在加法上是阿贝尔群，而其理想$I$则是$R$的正规加法子群，
因此$R/I$在加法上就是一个商群，其中元素就是加法上$I$的陪集。比如，任取$r\in R$，$r + I$
就是是$R/I$中的群元素。为清晰起见，描述$R/I$在加法定义如下。对任意的$r, s \in R$，
\[(r + I) + (s + I) = (r+s) + I\]
至此，我们没有引入任何新知识，已经对商环的元素和结构有了大致的了解。

接着，$R/I$要形成环，必须定义$R/I$群元素的乘法。无论乘法是什么，根据环的定义，该乘法必须是封闭的，且具有
结合律和分配律。在给出乘法定义之后，这些都需要证明。除此之外，特别强调的是，既然$R/I$的乘法
是对陪集的操作，千万要记得证明良定义属性，因为陪集的代表元不唯一。

\begin{lemma}{商环乘法}{qr_mul}
设$R$是环，$I$是$R$中的理想。商群$R/I$中元素的乘法定义为：
对任意的群元$r, s\in R$，
\[(r + I)(s +I) = rs + I\]
该乘法是一种良定义操作，且具有封闭性、结合律和对加法具有分配律。
\end{lemma}
要证明乘法是一种良定义操作，就是要证明乘法独立于陪集代表元的选择。即证明，如果$r + I = r' + I$，
$s + I = s' + I$，则$(r + I)(s + I) = (r'+ I)(s' + I) $。根据乘法定义，即要证$rs+ I = r's' + I$。
根据命题~\ref{pro:coset}，也就是要证明$r's' \in rs + I$。另外，同样根据命题~\ref{pro:coset}，
$r + I = r' + I$和$s + I = s' + I$分别代表$r' \in r + I$，$s' \in s + I$。
这就是证明思路。具体证明如下。
\begin{proof}
首先证明乘法是一种良定义操作，即假设$r' \in r + I$，$s' \in s + I$，证明$r's' \in (rs + I)$。
因为$r' \in r + I$，$s' \in s + I$，即存在$i_1, i_2 \in I$使得$r' = r + i_1$和$s' = s + i_2$，因此，
\[ r' s' = (r + i_1) (s + i_2) = rs + r i_2 + i_1 s + i_1 i_2\]
根据理想的吸收性，$r i_2 + i_1 s + i_1 i_2 \in I$，所以，$r's' \in rs + I $。
商环乘法的封闭性、结合律和分配律留作课后练习。\qed
\end{proof}

\begin{theorem}{商环}{Quo_ring}
设$R$是环，$I$是$R$中的理想。商群$R/I$在陪集加法与引理~\ref{lem:qr_mul}中定义的乘法上形成环，
称为$R$模$I$的商环，同样记为$R/I$。
\end{theorem}

\begin{example}
\label{exm:quo_ring}
任取$n \in \Z$，$n\Z = \langle n \rangle$是整数环$\Z$的主理想，则$\Z / n \Z$是商环，其中元素
刚好构成模$n$的完全剩余系。
\end{example}

%Canonical Homomorphism for Ring.
%20200712 morning
\begin{definition}{环的标准同态}{ring_can_homo}
		设$I$是环$R$的理想，定义环同态映射$\phi\colon R \to R / I$为：对任意$r \in R$，$\phi(r) = r + I$。
    并称该映射为环的\emph{标准同态}或者\emph{自然同态}，且$\Ker\; \phi = I$。 
\end{definition}

\begin{theorem}{第一同构定理}{ring_first_iso}
设$\psi\colon R \to S$是环同态，记$K= \Ker\; \psi$是$R$的理想。如果$\phi\colon R \to R/K$
是标准同态，则存在唯一同构$\eta\colon R/K \to \psi(R)$使得$\psi = \eta \phi$。
\end{theorem}

\begin{proof}
根据群论的第一同构定理（定理~\ref{thm:first_iso}），在$R$的加法群与$R$模$K$的加法商群之间，
存在唯一的良定义的群同构$\eta\colon R/K \to \psi(R)$。
该映射定义为，对任意的$r \in R$，有
\[\eta(r + K) = \psi(r)\]
要证明$\eta$是一种环同态，只需要证明，对任意的$r, s \in R$，有
\[\eta((r+K)(s+K))= \eta(r+K)\eta(s+K)\]
然而，这是容易的，因为
\begin{align*}
\eta((r+K)(s+K)) &= \eta(rs + K)\\
                                  &= \psi(rs)\\
                                  &= \psi(r)\psi(s)\\
                                  &= \eta(r+K)\eta(s+K)                                  
\end{align*}
\qed
\end{proof}

\begin{example}
\label{exm:Zn_mod_nZ}
任取$n \in \Z$，构造映射$\phi\colon \Z \to \Z_n$为，任取$a \in \Z$，$\phi(a) = a \bmod n$。可验证，这是
一个环同态映射。$\Ker\; \phi = n \Z$，因为所有$n$的倍数都映射为$0$。根据第一同构定理，
$\Z / n \Z \cong \Z_n$。
\end{example}

%%20200714
\section{域}
\label{sec:field}
\begin{definition}{域}{field}
如果一个带单位元的交换环$R$中的非$0$元素都存在唯一的乘法逆元，即
$\forall a \in R$且$a \neq 0$，则存在唯一的$a^{-1}\in R$使得$a \cdot a^{-1} = a^{-1}a =1$，则称这种代数结构
为域。
\end{definition}
一般用$\F$表示域\footnote{为什么群是$\mathbb{G}$，域是$\F$，但是环却不配写作$\R$？因为$\R$特指实数域。}，
而$\F^*$表示域中所有的非$0$元素。由定义可知，域$\F$在加法上是阿贝尔群，
同时$\F^*$在乘法上是
阿贝尔群。通常可以把域看为特殊的环。但是，这样说并不公平，
从域的发展史看，域的提出本身具有其特定的意义与内涵。
很少有人重点关注有限元素的环，但是有限元素的域被定义为\emph{有限域}（\emph{finite field}）
\index{有限域}\index{finite field}，
这是非常重要的概念，得到广泛的重视与研究，而且往往称
有限域为\emph{伽罗瓦域}（\emph{Galois field}）\index{伽罗瓦域}\index{Galois field}，
以此纪念该理论的先驱之一法国数学天才伽罗瓦（\'{E}variste Galois）。

本书不打算全面展开讨论域的相关理论，建议读者在此先把域看为环的特例，从讨论环与域的关系入手，
逐步熟悉域的属性。以此为基础再去学习域的理论则收获更大。

\begin{example}
\begin{enumerate} 
\item 在普通的加法与乘法上，$\mathbb{Q}, \mathbb{R}, \mathbb{C}$都是域。$\Z$不是域，因为乘法上$\Z$不是群。

\item 任意给定素数$p$，在模$p$的加法与乘法上$\Z_p$是域。因为$\Z_p$在加法上是阿贝尔群，而$\Z_p^*$在乘法上
也是阿贝尔群。对任意合数$n\in \Z$，$\Z_n$则不是域，因为$\Z_n - \{ 0\}$乘法上不成群。注意，因为域$\Z_p$使用很多，
为避免与加法群混淆，为此给出一种特定的标识符，记为$\F_p$。
\end{enumerate}
\end{example}

根据定义，域$\F$中的非$0$元素都存在唯一的乘法逆元，而这意味着$\F$中没有零因子，表述为以下命题。
\begin{proposition}{乘法逆元与零因子}{inverse_zero}
对任意的域$\F$，任取$a, b \in \F$，如果$ab = 0$则$a = 0$或者$b = 0$。即域中不存在零因子。
\end{proposition}

\begin{proof}
不妨设$a \neq 0$，否则证完。因为$a \neq 0$，则存在$a$的乘法逆元$a^{-1} \in \F$且$a^{-1} \neq 0$，使得$a a^{-1} =1$。
等式$ab = 0$两边乘上$a^{-1}$，根据命题~\ref{pro:ring0}，则$b = 0$。
\end{proof}
因此，所有的域都是整环。但是显然并非所有的
整环都是域，比如，$\Z$是整环，但并非域。以下结论则颇有点出人意料。

\begin{proposition}{整环与域}{int_domain_field}
每一个有限整环都是域。
\end{proposition}

\begin{proof}
证明的思路就是利用有限整环的性质，为每一个非$0$元素找到乘法逆元。
设$D$是一个有限整环，记$D^*$为环中所有非$0$元素的集合。对任意的$a \in D^*$，
构造映射$\lambda_a\colon D^* \to D^*$为$\lambda_a(d) = ad$，$\forall d \in D^*$。
首先，这确实是一个合理的映射，因为根据整环的性质，若$a \neq 0$和$d \neq 0$，则$ad \neq 0$。
其次，可证明$\lambda_a$是单射，因为对任意的$d_1, d_2 \in D^*$，如果：
\[ad_1 = \lambda_a(d_1) = \lambda_a(d_2) = a d_2\]
根据整环的消去律，则有$d_1 = d_2$。
然后，因为$D^*$是有限集且$\lambda_a$是从$D^*$到$D^*$的单射，所以$\lambda_a$必然是满射。
因此，必然存在某个$d\in D^*$使得$ad = 1$，又因为$D$是交换环，所以这个$d$就是$a$的乘法逆元。
结论：可为$D$中每一个非零元素都找到乘法逆元，所以$D$是一个域。
\end{proof}

\begin{remark}
由以上证明，读者能体会出为什么要求整环是交换环吗？
\end{remark}
在作者看来，以上证明的方法较为独特有趣，值得推荐给读者。如果只是为了证明，使用以下简单的构造法即可。注意到，
对任意非零元素$a \in D^*$，构造序列$a, a^2, a^3, \ldots$，必然存在两个正整数$i$和$j$使得$i > j$且$a^i = a^j$。
简单的计算可求出$a$的乘法逆元，细节留给读者自行完成。%%% 20230808添加，源自Gal21.

\begin{definition}{特征}{char}
环$R$的\emph{特征}（\emph{characteristic}）\index{特征}\index{characteristic}
定义为最小的正整数$n$使得对任意的$r\in R$，$\underbrace{r + r + \cdots + r}_{n\text{个}} = nr = 0 $。
如果不存在这样的$n$，则$R$的特征定义为$0$。
\end{definition}

\begin{example}
\label{exm:char}
\begin{enumerate} 
\item 环$\mathbb{Q}, \mathbb{R}, \mathbb{C}$的特征都是$0$。

\item 对任意素数$p$，域$\Z_p$的特征是$p$。因为$\Z_p$加法群的阶为$p$，
即对任意的$a \in \Z_p$，$pa = 0$。
\end{enumerate}
\end{example}

\begin{lemma}{}{one_char}
设$R$为带乘法单位元的环，若$1$在加法群的阶为$n$，则$R$的特征为$n$。
\end{lemma}

\begin{proof}
若$1$在加法群的阶为$n$，则$n$是最小正整数使得$n1 = 0$。那么，任取$r\in R$，有：
\[n r = n(1 r) = (n1)r = 0r = 0\]
即$n$是$R$的特征。
\end{proof}

\begin{proposition}{}{domain_char}
整环的特征或者为素数，或者为$0$。
\end{proposition}

\begin{proof}
设整环$D$的特征为$n$，且$n \neq 0$。
如果$n$不是素数，则$n = a b$，且$1 < a, b < n$。
根据引理~\ref{lem:one_char}，有
\[ 0 = n 1 = (ab)1 = (a1)(b1)\]
因为$D$是整环，$D$中无零因子，所以必然$a1 = 0$或者$b1=0$。
但是这都意味着$D$的特征小于$n$，矛盾。
\end{proof}

因为所有的域都是整环，所以域的特征也必然或者为素数，或者为$0$。
例~\ref{exm:char}已经展示了相关域的特征，其中域$\F_p$的特征与$\F_p$的阶都是$p$
有没有引起你的联想呢？域的特征与阶是什么样的关系呢？以下考虑有限域。

\begin{proposition}{有限域的特征}{ff_char}
阶为$n$的有限域$\F$的特征是一个素数$p$，且$p\mid n$。
\end{proposition}

\begin{proof}
因为有限域$\F$的阶为$n$，且$\F$是加法群，
所以对任意的$a\in \F$，有$na = 0$。
所以，$\F$的特征必然是素数$p$，且$p \mid n$。
\end{proof}

以下不加证明给出另一个重要结论。
\begin{proposition}{有限域的阶}{ff_order}
如果有限域$\F$的特征是素数$p$，则$\F$的阶是$p^n$，$n$是某个正整数。
进一步，对任意的素数$p$和正整数$n$，存在阶为$p^n$的有限域，并且
所有的$p^n$阶有限域都同构。
\end{proposition}

$p^n$阶的有限域也记为$\mathrm{GF}(p^n)$，$\mathrm{GF}$是 Galois field 的缩写。
这个结论最吸引计算机人的地方在于，因为$2$是素数，所以必然可以构造有$2^n$个元素的有限域。
如果说$\Z_n$是计算机系统最自然的代数结构，那么$\mathrm{GF}(2^n)$则是计算机系统的最理想化的代数结构。
\begin{example}
使用 SageMath 可轻松定义使用$\mathrm{GF}$域，请看以下简单例子。虽然简单，但是有些输出结果应该看不明白。
解释这些结果是这是接下来的任务，请先建立直观印象。
\begin{lstlisting}
\label{exm:GF_field}
sage: GF17 = GF(17)
sage: GF17.is_field()
True
sage: GF17.characteristic()  #域的特征
17
sage: GF17.order()    #域的阶
17
sage: GF17_2 = GF(17^2)
GF17.is_field()
True
sage: GF17_2.characteristic()
17
sage: GF17_2.order()
289
sage: GF17.random_element()   #输出一个随机域元素
11
sage: GF_17_2.random_element()
4*z2 + 3
\end{lstlisting}
\end{example}

是时候讨论域的同态、同构、同构定理、“商域”这些知识点了，然而如果视域为特殊的环，其实可讨论的并不多。
因为任何一个从域$\F_1$到域$\F_2$的域同态都只是从环$\F_1$到环$\F_2$的环同态。
所以，命题~\ref{pro:ring_homo}、定理~\ref{thm:Quo_ring}和定理~\ref{thm:ring_first_iso}等
对于域同样成立。特殊之处
在于，每一个域同态都是单射，或者是零同态。

在没有证明之前，仅仅知道结论，“域同态是单射”会让读者立即想到什么呢？
比如，考虑一个从域$\F_1$到域$\F_2$的域单射同态$\phi$，$\Ker\; \phi$会是什么样？
单射当然意味着$\Ker\; \phi$只有一个元素$0$。$\Ker\; \phi$是什么？域的理想！
所以，需要讨论任意域$\F$的理想所呈现的状态，有以下命题了。

\begin{proposition}{域的理想}{field_ideal}
任何一个域$\F$的理想只有$0$和自己本身$\F$。
\end{proposition}
\begin{proof}
首先，已知$0$和$\F$都是$\F$的理想。
设$I$是域$\F$的非$0$理想，则存在非零元素$a \in I$。因为$\F$是域，则存在$a$的乘法逆元$a^{-1} \in \F$。
根据理想的吸收性，$a^{-1}a = 1 \in I$。包含$1$的理想$I$等于$\F$，即$I = \F$。
\end{proof}

\begin{proposition}{域同态是单射}{homo_field}
任何一个域同态或者是单射或者是零同态。
\end{proposition}

\begin{proof}
域$\F_1$到域$\F_2$的域同态$\phi$是单射，当且仅当$\Ker\; \phi = \{0\} \subset \F_1$。
因为$\F_1$的理想只有$0$和$\F_1$本身，所以当$\Ker\; \phi = \{0\} $时$\phi$是单射，
而当$\Ker\; \phi = \F_1$时，$\phi$是零同态。
\end{proof}

既然域是特殊的环，有商环自然就有商域？但是，考虑到任何一个域$\F$的理想只有$0$和自己本身$\F$。
$\F$模$0$就是$\F$本身，而$\F$模$\F$就同构于$0$。也就是说，“商域”其实不太值得讨论。
值得思考的是，给定一个环$R$，什么样的商环$R/I$会是域？

虽然暂时还没有答案，但是环的第一同构定理已经给出了具体的实例。例~\ref{exm:Zn_mod_nZ}告诉我们，
任取$n \in \Z$，$\Z / n \Z \cong \Z_n$。任取素数$p \in \N$，立即有$\Z / p\Z \cong \Z_p$，而已知
$\Z_p$是域。所以，讨论理想$p\Z$的属性将会帮助我们找到正确答案。

\begin{definition}{极大理想与素理想}{max_pri}
设$R$是环，$M$是$R$的真子集且是$R$的理想，则称$M$是$R$的\emph{真理想}。
设$M$是$R$的真理想，
如果$M$不是$R$的任意真理想的真子集，则称
$M$是\emph{极大理想}（\emph{maximal ideal}）\index{极大理想}\index{maximal ideal}。
即如果$M$是$R$的极大理想，则对$R$的任意理想$I$，若$M \subset I$，则$I = R$。
设$P$是交换环$R$的真理想，如果对任意$ab \in P$，则或者$a \in P$，或者$b \in P$，
就称$P$为\emph{素理想}（\emph{prime ideal}）\index{素理想}\index{prime ideal}。
\end{definition}

\begin{example}
设$P = \{0, 2, 4, 6\}$为环$\Z_8$的理想，可验证$P$是极大理想，也是素理想。
任取素数$p$，则$p \Z$是$\Z$的素理想。
\end{example}

\begin{example}
素理想的“素”确实有“素数”的意味。请回顾第~\ref{chap:prime}章的定理~\ref{thm:euclid1}：
设$p$为素数，如果$p \mid ab$，则$p \mid a$或$p \mid b$。请体会素理想定义中要求与之类似之处：
对任意$ab \in P$，则$a \in P$或$b \in P$。令$P=p \Z$，$p$是任意素数，$a \in P$当且仅当$p \mid a$。
所以，从整数的角度上看，素理想确实是素数的倍数形成的理想。当然，理想不能仅停留于此，还需要进一步的抽象，
但是这个例子告诉我们，抽象代数的“抽象”并非凭空而出，往往源自于具体的实例。
\end{example}

\begin{theorem}{极大理想与域}{max_ideal}
设$R$是带单位元的交换环，$M$是$R$的理想，则$M$是$R$的极大理想，当且仅当$R/M$是域。
\end{theorem}

\begin{proof}
\begin{enumerate}
\item $\Rightarrow$. 由条件可知$R/M$是交换环，只需要证明$R/M$中所有非零元都有乘法逆元，则$R/M$是域。%即不包括0 + M
定义环同态$\phi\colon R \to R/M$为$\phi(r) = r + M$，$\forall r \in R$。任取$r \in R$且$r \notin M$，构造主理想$\langle r + M \rangle$。
因为$\phi$是环同态，所以$I = \phi^{-1}(\langle r + M \rangle)$是$R$中的理想。并且$M$是$I$的真理想，因为$M \subseteq I$，且存在$r\in I$但是$r\notin M$。
又因为$M$是极大理想，所以$I = R$。因此有：\[\phi(1) = 1 + M \in \langle r + M \rangle\]
即存在$s \in R$使得$1 + M = (r + M) (s + M)$。因此，$r + M$有乘法逆元。
%以上证明源自这里。
%https://ttjones.wordpress.com/2011/08/29/rm-is-a-field-iff-m-is-maximal/ 

\item $\Leftarrow$. 因为$R/M$是域，所以它至少包含两个元素$0 + M$和$1 + M$，即$M$是$R$的真理想。假设存在$R$的理想$I$，且$M$是$I$的真子集，
只需证$I = R$，则可证明$M$是$R$的极大理想。取$a \in I$且$a \notin M$，可知$a + M$是域$R/M$的非零元素。因此存在$b + M \in R/M$使得：
\[(a +M)(b +M) = ab +M = 1 + M \]
即存在$m \in M$使得$ ab + m = 1$，根据环的封闭性和理想的“吸收性”，可知$1 \in I$。所以，$I = R$。
\end{enumerate}
\end{proof}

\begin{example}
任取素数$p$，已知$p \Z$是$\Z$的素理想。因为$\Z / p\Z \cong \Z_p$，$\Z_p$是域，所以$p \Z$是$\Z$的极大理想，
\end{example}

\begin{theorem}{素理想与整环}{prime_ideal}
设$R$是带单位元的交换环，$P$是$R$的理想，则$P$是$R$的素理想，当且仅当$R/P$是整环。
\end{theorem}

\begin{proof}
\begin{enumerate}
\item $\Rightarrow$. 设$P$是素理想，若$R/P$中的两个元素使得：
\[(a + P) (b + P) = ab + P = 0 + P = P\]
可知，$ab \in P$。不失一般性，若$a \notin P$，则根据素理想的定义，有$b \in P$。因此，有$b + P = 0 + P$。即$R/P$是整环。

\item $\Leftarrow$. 设$P$是$R$的理想，且$R/P$是整环。假定$ab \in P$。则有：
\[(a + P) (b + P) = ab + P = 0 + P = P\]
根据整环属性，则或者$a + P = P$成立，或者$b + P = P$成立。这意味着，或者$a \in P$或者$b \in P$。说明$P$必须是素理想。
\end{enumerate}
\end{proof}

%%% 20230619晚改写
\begin{example}
已知$\Z$的理想都形如$n \Z$，$n$是整数。可知，仅当$n$为素数时，$\Z / n \Z \cong \Z_n$是整环，且是域。因此，$\Z$的素理想只能是$p \Z$，$p$是素数。
\end{example}

\begin{corollary}{极大理想与素理想}{max_prim_ideal}
交换环的每一个极大理想都是素理想。
\end{corollary}

\begin{proof}
容易。因为，所有的域都是整环。
\end{proof}

 %%2023年6月19日
 \section*{章注}
 %\section*{Chapter Note}
 \addcontentsline{toc}{section}{章注}
 本章的教学内容以群论为基础，在讲解环与域的知识点时，对照群论的同态、kernel、正规子群、商群和第一同构定理等概念，指出在环与域中这些概念的对应概念，
 让读者更好地把握整个知识链，帮助记忆与学习。同时，讲解过程也注重指出环、域、群之间的微妙区别。相信这种以群论为基础的教学只是一种选择，有不少
 抽象代数的教材从环这个概念开始教学。
 
%%%%%%%%%%%%%%%%%%%%%%%%%%%%%%%%%%%%%%%%%%%%
%%%课后习题
%%%%%%%%%%%%%%%%%%%%%%%%%%%%%%%%%%%%%%%%%%%%
 \begin{problemset}
   \item 如果环$R$带乘法单位元$1$，对任意$a\in R$，请证明$-a = (-1) a$。% (-1)a + a = ((-1) + 1) a = 0 a = 0 。
   %源自AATA
 
  \item 如果任取环$R$中的元素$x$都满足$x^2 = x$，请证明环$R$是交换环。
  
   %%% From Gal 21 第12章41题，他的答案不好。
   %%% 答案：因为 a*b = a + ... + a (b个a相加)，加法封闭，所以必然有乘法封闭。
  \item 请解释为什么$\Z_n$在加法上的子群都是$\Z_n$的子环。
  
  %%% From Gal 21 第12章12题。容易。
  \item 设$a, b, c$是交换环$R$中的元素，且假定$a$有乘法逆元。请证明$b$整除$c$，当且仅当$ab$整除$c$.
  
  %%% From Gal 21 第13章课后习题第五题。
  \item 证明$\Z_n$中的所有非零元素，或者是一个零因子，或者该元素存在乘法逆元。
  
  \item 记$\Z[\sqrt{2}] = \{a + b\sqrt{2}: a, b \in \Z \}$，请证明$\Z[\sqrt{2}] $是环，且是整环。
  
  \item 记$\Z[i] = \{a + bi: a, b \in \Z \}$，请证明$\Z[i]$是整环。%AATA 
 
  \item 如果$I$和$J$都是交换环$R$的理想，则$I+J = \{ i + j : i \in I, j \in J\}$也是$R$的理想。 
  
  \item 证明命题~\ref{pro:ring_homo}。
  
  \item 任取$n \in \Z$，构造映射$\phi\colon \Z \to \Z_n$为，任取$a \in \Z$，$\phi(a) = a \bmod n$。请验证
  $\phi$是一个环同态映射。
  
  \item 请验证例子~\ref{exm:homo_onto}的同态实例，并回答，为什么命题~\ref{pro:ring_homo}要求$\phi$是满射才会有$\phi(1) = 1'$？
  
  \item 请补充完成引理~\ref{lem:qr_mul}的证明。
  
  \item 设$R$是环，$I$和$J$是$R$中的理想，并且$I + J = R$。请证明存在以下环同构：
  \[R/(I \cap J) \cong R/I \times R/J\]
  提示：使用定理~\ref{thm:crt_ring}和环的第一同构定理。
  
  \item 证明环$2\Z$不与环$3\Z$同构。%https://math.stackexchange.com/questions/1277844/show-that-2z-and-3z-are-not-isomorphic-question-on-proof
  %6k = 9k^2 iff k = 0.
  
  \item 证明环$\R$不与环$\mathbb{C}$同构。%反证法，如果存在同构f使得f(i) = a，因为f(1) = 1, 那么f(-1)  = -1 （0 = f(0) = f(1 + (-1)) = f(1) + f(-1)，所以有f(-1) = -f(1) = -1）。f(i^2) = -1 = a^2。矛盾！20250501补充。
  %https://math.stackexchange.com/questions/101765/prove-that-mathbbc-and-mathbbr-are-not-isomorphic-as-rings
  
  \item 记$\Q(\sqrt{d})  = \{a + b \sqrt{d}: a, b \in \Q\}$，$d$是任意有理数，且是非平方数。证明或证伪：整环$\Q(\sqrt{2}) $与整环$\Q(\sqrt{3})$同构。 
  %https://yutsumura.com/two-quadratic-fields-are-not-isomorphic/
  %https://math.stackexchange.com/questions/9188/is-mathbbq-sqrt2-cong-mathbbq-sqrt3
  %https://math.stackexchange.com/questions/302514/are-mathbbr-and-mathbbc-isomorphic-as-additive-groups
  
 \item 设$p$是有限域$\mathbb{F}_q$的特征，其中$q=p^n$，$n$是一个正整数。请证明当$p=2$时，$\mathbb{F}$中的任意元素都是一个平方数，即对任意$a\in \mathbb{F}_q$，存在$b\in \mathbb{F}_q$，使得$a = b^2$。如果$p$是一个奇数，则$\mathbb{F}$中一半的非零元素是平方数。%%% Hil86-FCCT
 
 % 可能来自FCAA，已经写了答案。
 \item  证明：有限域中任意一个元素可以表达为该有限域中两个元素的平方和。即对任意有限域$\mathbb{F}_q$ 的任意元素$c\in \mathbb{F}_q$，存在$a, b \in \mathbb{F}_q$使得$c = a^2 + b^2$。
  \end{problemset}
